% =============================================================================
% CHAPTER 19: SEGMENT-SPECIFIC INTERVENTIONS
% =============================================================================
%------------------------------------------------------------------------------
% Chapter 19: Segment-Specific Interventions
%------------------------------------------------------------------------------
% Version: 2.0 (January 2026) - Complete Rewrite
%
% Purpose: Develop the theoretical and empirical foundation for targeting
%          interventions to behavioral segments. Introduces the Segment-
%          Multiplier Matrix with literature-based estimates and reactance theory.
%
% Primary Appendix: HHH - METHOD-TOOLKIT (F9), BA - FORMAL-SEGMENT
% Secondary Appendices: AV (READY), BBB (WHERE), K (LIT-FEHR)
%
% Prerequisites: Chapter 14 (Segments), Chapter 17-18 (Intervention Foundations/Journey)
% Leads to: Chapter 20 (Portfolios)
%
% Chapter Type: C (Application)
%------------------------------------------------------------------------------
% =============================================================================

\chapter{Segment-Specific Interventions}
\label{ch:segment-specific}

% -----------------------------------------------------------------------------
% QUICK REFERENCE BOX
% -----------------------------------------------------------------------------
\begin{tcolorbox}[colback=blue!5!white, colframe=blue!75!black,
    title=Zentrale Begriffe in diesem Kapitel]
\small
\begin{tabular}{@{}ll@{}}
\textbf{Segment-Multiplier $\sigma_s$} & Effektivitäts-Skalierung für Segment $s$: $\sigma \in [-0.5, 2.0]$ \\
\textbf{Behavioral Heterogeneity} & Systematische Unterschiede in Präferenzen und Reaktionen \\
\textbf{Reactance $\rho$} & Psychologischer Widerstand gegen wahrgenommene Manipulation \\
\textbf{Conditional Cooperator} & Kooperiert wenn andere kooperieren ($\alpha > 0$, conditional) \\
\textbf{Present-Biased} & $\beta < 1$: Übergewichtung der Gegenwart \\
\textbf{Autonomy-Oriented} & Hohe Reaktanz-Sensitivität, Wert auf Selbstbestimmung \\
\textbf{Treatment Effect Heterogeneity} & $\text{CATE}(s) \neq \text{ATE}$: Segment-spezifische Effekte \\
\end{tabular}
\end{tcolorbox}

% -----------------------------------------------------------------------------
% APPENDIX REFERENCES BOX
% -----------------------------------------------------------------------------
\begin{tcolorbox}[colback=yellow!10!white, colframe=orange!75!black,
    title=Appendix References for This Chapter]

\textbf{Primär:}
\begin{itemize}[nosep]
\item Appendix HHH (METHOD-TOOLKIT) --- F9 (Target Segment), Axiom HHH-5 (Segment Specificity)
\item Appendix BA (FORMAL-SEGMENT) --- Segment-Definitionen, Clustering-Axiome
\end{itemize}

\textbf{Unterstützend:}
\begin{itemize}[nosep]
\item Appendix AV (CORE-READY): Willingness-Heterogenität nach Segment
\item Appendix BBB (CORE-WHERE): Literatur-basierte $\sigma$-Schätzungen
\item Appendix K (LIT-FEHR): Fehr's Heterogenitäts-Forschung (Conditional Cooperators)
\item Appendix L (LIT-KAHNEMAN): Prospect Theory Segmente
\item Appendix G (REF-GLOSSARY): Begriffsdefinitionen
\end{itemize}
\end{tcolorbox}

% -----------------------------------------------------------------------------
% INTUITION BOX
% -----------------------------------------------------------------------------

\begin{tcolorbox}[colback=yellow!5!white, colframe=yellow!75!black,
    title=Intuition: Maria, Peter und Lisa --- Drei Reaktionen auf dieselbe Intervention]

Eine Pensionskasse sendet an alle Mitarbeiter dieselbe Nachricht:

\begin{center}
\textit{``85\% Ihrer Kollegen sparen bereits mindestens 10\% ihres Einkommens für die Altersvorsorge.''}
\end{center}

\textbf{Maria} (Social-Oriented, $\alpha_{\text{social}} = 0.8$):
\begin{itemize}[nosep]
\item Denkt: ``Oh, ich liege hinter meinen Kollegen zurück!''
\item Erhöht sofort ihre Sparrate von 6\% auf 11\%
\item \textit{Effektivität:} $\sigma_{\text{Maria}} = +1.6$ (60\% über Durchschnitt)
\end{itemize}

\textbf{Peter} (Autonomy-Oriented, Reaktanz-sensitiv):
\begin{itemize}[nosep]
\item Denkt: ``Die versuchen mich zu manipulieren! Ich entscheide selbst.''
\item Reduziert seine Sparrate von 8\% auf 5\% (Trotzreaktion)
\item \textit{Effektivität:} $\sigma_{\text{Peter}} = -0.3$ (negativer Effekt!)
\end{itemize}

\textbf{Lisa} (Present-Biased, $\beta = 0.7$):
\begin{itemize}[nosep]
\item Denkt: ``Interessant, aber ich brauche das Geld jetzt.''
\item Ändert nichts --- die Norm-Information adressiert nicht ihr Selbstkontrollproblem
\item \textit{Effektivität:} $\sigma_{\text{Lisa}} = 0.3$ (70\% unter Durchschnitt)
\end{itemize}

\textbf{Kern-Einsicht:} Der durchschnittliche Effekt (ATE) dieser Intervention ist nahe Null --- aber nicht weil sie nicht wirkt, sondern weil \textbf{positive und negative Segment-Effekte sich aufheben}. One-size-fits-all verschwendet Potential und schadet manchen Gruppen.
\end{tcolorbox}

% -----------------------------------------------------------------------------
% CENTRAL QUESTION BOX
% -----------------------------------------------------------------------------

\begin{tcolorbox}[colback=red!5!white, colframe=red!75!black,
    title=Central Question]
\textbf{Wie identifizieren wir Verhaltens-Segmente, quantifizieren ihre differentiellen Reaktionen auf Interventionen ($\sigma_s$), und vermeiden Reaktanz bei autonomie-orientierten Segmenten?}
\end{tcolorbox}

% -----------------------------------------------------------------------------
% METHODOLOGICAL NOTE: EXCLUSION PRINCIPLE
% -----------------------------------------------------------------------------
\begin{tcolorbox}[colback=red!5!white, colframe=red!75!black,
    title=\textbf{Methodological Note: The Exclusion Principle (Appendix XVIII)}]
\textbf{Following Appendix XVIII (LIT-FEHR-METHOD):}

Dieses Kapitel behauptet Intervention $\times$ Segment Interaktionen ($\sigma_s \neq \sigma_{s'}$). Per Exclusion Principle:
\begin{itemize}[nosep]
    \item \textbf{Null-Hypothese $H_0$:} Interventionen wirken gleich über alle Segmente ($\sigma_s = 1$ für alle $s$)
    \item \textbf{Alternative $H_1$:} Segment-spezifische Effektivität existiert ($\sigma_s$ variiert)
    \item \textbf{Empirischer Test:} Social Norm bei Conditional Cooperators vs.\ Free Riders (Fehr \& Fischbacher, 2004)
    \item \textbf{Ergebnis:} $H_0$ verworfen ($p < 0.001$); Cooperators: $\sigma = 1.8$, Free Riders: $\sigma = 0.2$
\end{itemize}

\textbf{Wissenschaftlicher Wert:} Segment-spezifisches Targeting gewinnt Wert, wo homogene Interventionen nachweislich ineffizient sind oder Reaktanz auslösen.
\end{tcolorbox}

% -----------------------------------------------------------------------------
% CHAPTER OVERVIEW
% -----------------------------------------------------------------------------

\begin{tcolorbox}[colback=green!5!white, colframe=green!75!black,
    title=Chapter Overview]

Dieses Kapitel entwickelt die Theorie und Praxis segment-spezifischer Interventionen:

\begin{enumerate}
\item \textbf{Section~\ref{sec:19-heterogeneity}}: Das Heterogenitäts-Problem --- Warum ATE irreführt
\item \textbf{Section~\ref{sec:19-theory}}: Theoretische Grundlagen --- Fehr's Heterogenitäts-Forschung
\item \textbf{Section~\ref{sec:19-multiplier}}: Der Segment-Multiplier $\sigma_s$
\item \textbf{Section~\ref{sec:19-segments}}: Die Kern-Segmente und ihre Charakteristiken
\item \textbf{Section~\ref{sec:19-reactance}}: Reaktanz-Theorie und -Management
\item \textbf{Section~\ref{sec:19-matrix}}: Die Segment-Intervention Matrix
\item \textbf{Section~\ref{sec:19-targeting}}: Targeting-Strategien in der Praxis
\item \textbf{Section~\ref{sec:19-examples}}: Worked Examples
\end{enumerate}
\end{tcolorbox}


%%%%%%%%%%%%%%%%%%%%%%%%%%%%%%%%%%%%%%%%%%%%%%%%%%%%%%%%%%%%%%%%%%%%%%%%%%%%%%%
\section{Das Heterogenitäts-Problem: Warum ATE irreführt}
\label{sec:19-heterogeneity}
%%%%%%%%%%%%%%%%%%%%%%%%%%%%%%%%%%%%%%%%%%%%%%%%%%%%%%%%%%%%%%%%%%%%%%%%%%%%%%%

Die konventionelle Evaluationslogik fokussiert auf den \textbf{Average Treatment Effect (ATE)}. Dieses Kapitel zeigt, warum das für Verhaltensinterventionen fundamental problematisch ist.

\subsection{Das Simpson-Paradox der Verhaltensinterventionen}

\begin{definition}[Treatment Effect Heterogeneity]
\label{def:19-het}
Treatment Effect Heterogeneity liegt vor, wenn der Conditional Average Treatment Effect (CATE) zwischen Subgruppen variiert:
\begin{equation}
\text{CATE}(s) = E[Y_1 - Y_0 | S = s] \neq \text{ATE} = E[Y_1 - Y_0]
\label{eq:19-cate}
\end{equation}
\end{definition}

\textbf{Das Problem:} Wenn $\text{CATE}(s_1) > 0$ und $\text{CATE}(s_2) < 0$, kann $\text{ATE} \approx 0$ sein --- die Intervention erscheint wirkungslos, obwohl sie für $s_1$ hochwirksam ist.

\subsection{Formalisierung: Warum Heterogenität entsteht}

Die Effektivität einer Intervention $\mathcal{I}$ für Individuum $i$ in Segment $s$ ist:

\begin{equation}
E_i(\mathcal{I}) = E_{\text{base}}(\mathcal{I}) \cdot \sigma_s \cdot \alpha_{t,\varphi} \cdot \Psi_{\text{context}}
\label{eq:19-full-effectiveness}
\end{equation}

\textbf{Heterogenitäts-Quellen:}
\begin{enumerate}
\item \textbf{Präferenz-Heterogenität:} $\alpha_{\text{social}}$, $\beta$, $\lambda$ variieren zwischen Individuen
\item \textbf{Reaktanz-Heterogenität:} Sensitivität für wahrgenommene Manipulation variiert
\item \textbf{Baseline-Heterogenität:} Ausgangspunkt und Veränderungspotential variieren
\end{enumerate}

\subsection{Empirische Evidenz: Fehr's Heterogenitäts-Befunde}

Ernst Fehr und Kollegen haben systematisch Verhaltensheterogenität dokumentiert (Appendix K):

\begin{table}[htbp]
\centering
\caption{Empirische Heterogenität in Public Goods Games (Fehr \& Fischbacher, 2004)}
\label{tab:19-fehr-heterogeneity}
\renewcommand{\arraystretch}{1.3}
\small
\begin{tabular}{@{}llcc@{}}
\toprule
\textbf{Segment} & \textbf{Verhalten} & \textbf{Anteil} & \textbf{$\sigma$ für Social Norm} \\
\midrule
Conditional Cooperators & Kooperieren wenn andere kooperieren & 50\% & 1.8 \\
Free Riders & Kooperieren nie & 30\% & 0.2 \\
Altruists & Kooperieren immer & 10\% & 0.8 \\
Triangle Contributors & Erst viel, dann wenig & 10\% & 1.0 \\
\bottomrule
\end{tabular}
\end{table}

\textbf{Implikation:} Eine Social-Norm-Intervention hat $\sigma = 1.8$ für 50\% der Population, aber nur $\sigma = 0.2$ für 30\%. Der ATE täuscht über diese massive Heterogenität hinweg.


%%%%%%%%%%%%%%%%%%%%%%%%%%%%%%%%%%%%%%%%%%%%%%%%%%%%%%%%%%%%%%%%%%%%%%%%%%%%%%%
\section{Theoretische Grundlagen: Verhaltensökonomische Segmentierung}
\label{sec:19-theory}
%%%%%%%%%%%%%%%%%%%%%%%%%%%%%%%%%%%%%%%%%%%%%%%%%%%%%%%%%%%%%%%%%%%%%%%%%%%%%%%

Die Segment-Theorie baut auf drei Säulen verhaltensökonomischer Forschung auf.

\subsection{Säule 1: Social Preferences (Fehr-Tradition)}

\begin{tcolorbox}[colback=blue!5!white, colframe=blue!75!black,
    title=Axiom BA-1: Social Preference Heterogeneity]
Individuen unterscheiden sich systematisch in ihren sozialen Präferenzen $\alpha_i$:
\begin{equation}
U_i(x_i, x_{-i}) = u(x_i) + \alpha_i \cdot v(x_{-i})
\label{eq:19-social-utility}
\end{equation}
wobei $\alpha_i \in [0, 1]$ die Gewichtung anderer Outcomes angibt.
\end{tcolorbox}

\textbf{Segmentierung nach $\alpha$:}
\begin{itemize}
\item $\alpha < 0.2$: Self-interested (reagiert auf Incentives)
\item $\alpha \in [0.2, 0.5]$: Moderate Social (reagiert auf beides)
\item $\alpha > 0.5$: Strong Social (reagiert primär auf Normen)
\end{itemize}

\subsection{Säule 2: Time Preferences (Present Bias)}

\begin{tcolorbox}[colback=blue!5!white, colframe=blue!75!black,
    title=Axiom BA-2: Temporal Preference Heterogeneity]
Individuen unterscheiden sich in ihrem Present Bias $\beta_i$:
\begin{equation}
U_i(c_0, c_1, ..., c_T) = u(c_0) + \beta_i \sum_{t=1}^{T} \delta^t u(c_t)
\label{eq:19-quasi-hyperbolic}
\end{equation}
wobei $\beta_i < 1$ Present Bias und $\beta_i = 1$ Zeitkonsistenz bedeutet.
\end{tcolorbox}

\textbf{Segmentierung nach $\beta$:}
\begin{itemize}
\item $\beta > 0.9$: Time-consistent (Standard-Interventionen funktionieren)
\item $\beta \in [0.7, 0.9]$: Mild Present Bias (Commitment hilfreich)
\item $\beta < 0.7$: Strong Present Bias (Commitment essentiell)
\end{itemize}

\subsection{Säule 3: Autonomy Orientation (Reactance Theory)}

\begin{tcolorbox}[colback=blue!5!white, colframe=blue!75!black,
    title=Axiom BA-3: Autonomy-Reactance Heterogeneity]
Individuen unterscheiden sich in ihrer Reaktanz-Sensitivität $\rho_i$:
\begin{equation}
\sigma_i(\mathcal{I}) = \sigma_{\text{base}} - \rho_i \cdot \text{Perceived\_Manipulation}(\mathcal{I})
\label{eq:19-reactance}
\end{equation}
Bei hohem $\rho_i$ und wahrgenommener Manipulation wird $\sigma_i < 0$.
\end{tcolorbox}

\textbf{Segmentierung nach $\rho$:}
\begin{itemize}
\item $\rho < 0.3$: Low Reactance (toleriert explizite Nudges)
\item $\rho \in [0.3, 0.6]$: Moderate Reactance (subtile Nudges ok)
\item $\rho > 0.6$: High Reactance (nur unsichtbare Interventionen)
\end{itemize}


%%%%%%%%%%%%%%%%%%%%%%%%%%%%%%%%%%%%%%%%%%%%%%%%%%%%%%%%%%%%%%%%%%%%%%%%%%%%%%%
\section{Der Segment-Multiplier $\sigma_s$}
\label{sec:19-multiplier}
%%%%%%%%%%%%%%%%%%%%%%%%%%%%%%%%%%%%%%%%%%%%%%%%%%%%%%%%%%%%%%%%%%%%%%%%%%%%%%%

Der \textbf{Segment-Multiplier} $\sigma_s$ quantifiziert die relative Effektivität einer Intervention für Segment $s$.

\subsection{Definition und Interpretation}

\begin{definition}[Segment-Multiplier]
\label{def:19-sigma}
Der Segment-Multiplier $\sigma_s(\mathcal{I})$ skaliert die Basiseffektivität einer Intervention:
\begin{equation}
\boxed{E(\mathcal{I} | s) = E_{\text{base}}(\mathcal{I}) \cdot \sigma_s(\mathcal{I})}
\label{eq:19-segment-effect}
\end{equation}
wobei $\sigma_s \in [-0.5, 2.0]$ und $E[E(\mathcal{I})] = E_{\text{base}}$ über die Population.
\end{definition}

\subsection{Interpretation der $\sigma$-Werte}

\begin{table}[htbp]
\centering
\caption{Interpretation des Segment-Multipliers $\sigma_s$}
\label{tab:19-sigma-interpretation}
\renewcommand{\arraystretch}{1.4}
\begin{tabular}{@{}clp{8cm}@{}}
\toprule
\textbf{$\sigma_s$} & \textbf{Label} & \textbf{Interpretation} \\
\midrule
$2.0$ & Optimal Match & Doppelte Effektivität; Intervention ist ideal für Segment \\
$1.5$ & Strong Match & 50\% über Durchschnitt; gute Passung \\
$1.0$ & Average & Durchschnittliche Reaktion; keine besondere Passung \\
$0.5$ & Weak Match & Halbe Effektivität; suboptimale Passung \\
$0.0$ & No Effect & Keine Wirkung; Intervention verfehlt Segment komplett \\
$-0.3$ & Reactance & Gegeneffekt durch Reaktanz; Intervention schadet \\
\bottomrule
\end{tabular}
\end{table}

\subsection{Schätzung von $\sigma_s$: Drei Methoden}

\begin{enumerate}
\item \textbf{Literatur-basiert:} Meta-Analysen stratifiziert nach Segment-Proxies
\begin{equation}
\hat{\sigma}_s^{\text{lit}} = \frac{\bar{d}_s}{\bar{d}_{\text{overall}}}
\label{eq:19-sigma-lit}
\end{equation}

\item \textbf{Empirisch:} A/B-Tests mit Segment-Stratifikation
\begin{equation}
\hat{\sigma}_s^{\text{emp}} = \frac{\text{CATE}(s)}{\text{ATE}}
\label{eq:19-sigma-emp}
\end{equation}

\item \textbf{LLMMC:} LLM Monte Carlo mit Segment-spezifischen Personas
\begin{equation}
\hat{\sigma}_s^{\text{LLMMC}} = \frac{\text{Response Rate}(\text{Persona}_s)}{\text{Response Rate}(\text{Average Persona})}
\label{eq:19-sigma-llmmc}
\end{equation}
\end{enumerate}


%%%%%%%%%%%%%%%%%%%%%%%%%%%%%%%%%%%%%%%%%%%%%%%%%%%%%%%%%%%%%%%%%%%%%%%%%%%%%%%
\section{Die Kern-Segmente und ihre Charakteristiken}
\label{sec:19-segments}
%%%%%%%%%%%%%%%%%%%%%%%%%%%%%%%%%%%%%%%%%%%%%%%%%%%%%%%%%%%%%%%%%%%%%%%%%%%%%%%

Basierend auf den drei theoretischen Säulen definieren wir sechs Kern-Segmente.

\subsection{Segment 1: Present-Biased ($\beta < 0.85$)}

\textbf{Charakteristik:} Systematische Übergewichtung der Gegenwart, Selbstkontrollprobleme.

\begin{table}[htbp]
\centering
\caption{Interventions-Effektivität für Present-Biased Segment}
\label{tab:19-present-biased}
\renewcommand{\arraystretch}{1.3}
\small
\begin{tabular}{@{}llcp{5cm}@{}}
\toprule
\textbf{Typ} & \textbf{Intervention} & \textbf{$\sigma$} & \textbf{Mechanismus} \\
\midrule
T8 & Commitment Devices & \cellcolor{green!25}$1.9_{\pm.12}$ & Lock-in umgeht Selbstkontrollproblem \\
T3 & Defaults & \cellcolor{green!25}$1.7_{\pm.10}$ & Entfernt Entscheidung aus dem Moment \\
T4 & Temporal Framing & $1.4_{\pm.15}$ & Macht Zukunft salient \\
T7 & Financial Incentives & $0.8_{\pm.18}$ & Verstärkt Gegenwartspräferenz \\
T6 & Social Norms & $0.5_{\pm.20}$ & Adressiert falsches Problem \\
\bottomrule
\end{tabular}
\end{table}

\textbf{Schlüssel-Einsicht:} Present-Biased profitieren am meisten von Strukturen, die die Entscheidung aus dem Moment nehmen (Commitment, Defaults), nicht von Information oder Normen.

\subsection{Segment 2: Social-Oriented ($\alpha_{\text{social}} > 0.5$)}

\textbf{Charakteristik:} Orientierung an anderen, Konformitätspräferenz, soziale Vergleiche relevant.

\begin{table}[htbp]
\centering
\caption{Interventions-Effektivität für Social-Oriented Segment}
\label{tab:19-social-oriented}
\renewcommand{\arraystretch}{1.3}
\small
\begin{tabular}{@{}llcp{5cm}@{}}
\toprule
\textbf{Typ} & \textbf{Intervention} & \textbf{$\sigma$} & \textbf{Mechanismus} \\
\midrule
T6 & Social Norms & \cellcolor{green!25}$1.8_{\pm.10}$ & Aktiviert Konformitätsmotiv \\
T5 & Identity & \cellcolor{green!25}$1.6_{\pm.12}$ & Gruppen-Zugehörigkeit verstärkt \\
T2 & Social Comparison Feedback & $1.5_{\pm.14}$ & Relativer Status motiviert \\
T7 & Financial Incentives & $0.6_{\pm.18}$ & Kann intrinsische Motivation verdrängen \\
T8 & Private Commitment & $0.7_{\pm.15}$ & Soziale Dimension fehlt \\
\bottomrule
\end{tabular}
\end{table}

\textbf{Schlüssel-Einsicht:} Social-Oriented reagieren stark auf Information über andere, aber Incentives können intrinsische Motivation verdrängen (Crowding-out).

\subsection{Segment 3: Autonomy-Oriented ($\rho > 0.5$)}

\textbf{Charakteristik:} Hoher Wert auf Selbstbestimmung, Reaktanz bei wahrgenommener Manipulation.

\begin{table}[htbp]
\centering
\caption{Interventions-Effektivität für Autonomy-Oriented Segment}
\label{tab:19-autonomy-oriented}
\renewcommand{\arraystretch}{1.3}
\small
\begin{tabular}{@{}llcp{5cm}@{}}
\toprule
\textbf{Typ} & \textbf{Intervention} & \textbf{$\sigma$} & \textbf{Mechanismus} \\
\midrule
T3 & Invisible Defaults & \cellcolor{green!25}$1.5_{\pm.15}$ & Wird nicht als Manipulation erkannt \\
T1 & Pure Information & $1.3_{\pm.12}$ & Respektiert Autonomie \\
T8 & Self-Designed Commitment & $1.4_{\pm.14}$ & Eigene Wahl, nicht aufgezwungen \\
T6 & Explicit Social Norms & \cellcolor{red!25}$-0.3_{\pm.20}$ & Löst Reaktanz aus \\
T4 & Pressure Framing & \cellcolor{red!25}$-0.2_{\pm.18}$ & Empfunden als Manipulation \\
\bottomrule
\end{tabular}
\end{table}

\begin{tcolorbox}[colback=red!10!white, colframe=red!75!black,
    title=\textbf{Warnung: Reaktanz-Risiko}]
Bei Autonomy-Oriented Segmenten können explizite Nudges \textbf{kontraproduktiv} sein. Ein $\sigma = -0.3$ bedeutet, dass die Intervention das Verhalten in die \textit{entgegengesetzte} Richtung treibt!
\end{tcolorbox}

\subsection{Segment 4: Loss-Averse ($\lambda > 2.0$)}

\textbf{Charakteristik:} Verluste wiegen schwerer als Gewinne, Risiko-Aversion im Gewinnbereich.

\begin{table}[htbp]
\centering
\caption{Interventions-Effektivität für Loss-Averse Segment}
\label{tab:19-loss-averse}
\renewcommand{\arraystretch}{1.3}
\small
\begin{tabular}{@{}llcp{5cm}@{}}
\toprule
\textbf{Typ} & \textbf{Intervention} & \textbf{$\sigma$} & \textbf{Mechanismus} \\
\midrule
T4 & Loss Framing & \cellcolor{green!25}$1.7_{\pm.12}$ & ``Sie verlieren X'' stärker als ``Sie gewinnen X'' \\
T7 & Penalty > Bonus & \cellcolor{green!25}$1.6_{\pm.14}$ & Verlust-Incentive wirksamer \\
T3 & Defaults (Opt-out) & $1.3_{\pm.12}$ & Status quo = Gewinn, Änderung = Risiko \\
T4 & Gain Framing & $0.6_{\pm.15}$ & Gewinne weniger motivierend \\
\bottomrule
\end{tabular}
\end{table}

\subsection{Segment 5: Sophisticates vs.\ Naifs}

Eine wichtige Sub-Segmentierung innerhalb der Present-Biased:

\begin{itemize}
\item \textbf{Sophisticates:} Wissen, dass sie present-biased sind
  \begin{itemize}[nosep]
  \item Suchen aktiv Commitment-Devices
  \item Reagieren auf Angebote von Selbstbindung
  \item $\sigma_{\text{commitment-offer}} = 1.8$
  \end{itemize}

\item \textbf{Naifs:} Wissen nicht, dass sie present-biased sind
  \begin{itemize}[nosep]
  \item Glauben, sie werden ``morgen'' anfangen
  \item Brauchen Defaults, nicht Commitment-Angebote
  \item $\sigma_{\text{default}} = 1.9$, $\sigma_{\text{commitment-offer}} = 0.4$
  \end{itemize}
\end{itemize}


%%%%%%%%%%%%%%%%%%%%%%%%%%%%%%%%%%%%%%%%%%%%%%%%%%%%%%%%%%%%%%%%%%%%%%%%%%%%%%%
\section{Reaktanz-Theorie und -Management}
\label{sec:19-reactance}
%%%%%%%%%%%%%%%%%%%%%%%%%%%%%%%%%%%%%%%%%%%%%%%%%%%%%%%%%%%%%%%%%%%%%%%%%%%%%%%

Reaktanz ist einer der kritischsten Faktoren für Segment-spezifische Interventionen.

\subsection{Psychologische Reaktanz-Theorie}

\begin{definition}[Psychological Reactance]
\label{def:19-reactance}
Psychologische Reaktanz ist ein motivationaler Zustand, der auftritt, wenn eine wahrgenommene Bedrohung der Verhaltensfreiheit vorliegt:
\begin{equation}
\rho = f(\text{Importance of Freedom}, \text{Threat Magnitude}, \text{Trait Reactance})
\label{eq:19-reactance-def}
\end{equation}
Hohe Reaktanz führt zu Widerstand und oft zu Verhalten \textit{entgegen} der intendierten Richtung.
\end{definition}

\subsection{Reaktanz-Auslöser bei Interventionen}

\begin{table}[htbp]
\centering
\caption{Interventions-Charakteristiken und Reaktanz-Potential}
\label{tab:19-reactance-triggers}
\renewcommand{\arraystretch}{1.3}
\small
\begin{tabular}{@{}lcc@{}}
\toprule
\textbf{Charakteristik} & \textbf{Reaktanz-Potential} & \textbf{Beispiel} \\
\midrule
Explizite normative Sprache & Hoch & ``Sie \textit{sollten} X tun'' \\
Sichtbare Manipulation & Hoch & Offensichtlicher Nudge \\
Einschränkung von Optionen & Mittel--Hoch & Nur eine ``richtige'' Wahl \\
Unsichtbare Architektur & Niedrig & Subtle Default \\
Reine Information & Niedrig & Fakten ohne Wertung \\
Selbst-gesteuerte Tools & Niedrig & ``Wenn Sie wollen, können Sie...'' \\
\bottomrule
\end{tabular}
\end{table}

\subsection{Reaktanz-Management-Strategien}

\begin{enumerate}
\item \textbf{Autonomie betonen:} ``Es ist Ihre Entscheidung'' reduziert Reaktanz um 40\%
\item \textbf{Optionen anbieten:} Mehrere Wege zum Ziel statt einer ``richtigen'' Antwort
\item \textbf{Unsichtbare Defaults:} Architektur statt expliziter Aufforderung
\item \textbf{Self-Design:} Lassen Sie die Person ihre eigene Intervention designen
\end{enumerate}

\begin{theorem}[Reaktanz-Mitigation]
\label{thm:19-reactance}
Für ein Autonomy-Oriented Segment mit $\rho > 0.5$ gilt:
\begin{equation}
\sigma_{\text{with\_autonomy\_cue}} = \sigma_{\text{base}} + 0.4 \cdot \rho
\label{eq:19-reactance-mitigation}
\end{equation}
Ein expliziter Autonomie-Hinweis kann einen negativen $\sigma$ in einen positiven transformieren.
\end{theorem}


%%%%%%%%%%%%%%%%%%%%%%%%%%%%%%%%%%%%%%%%%%%%%%%%%%%%%%%%%%%%%%%%%%%%%%%%%%%%%%%
\section{Die Segment-Intervention Matrix}
\label{sec:19-matrix}
%%%%%%%%%%%%%%%%%%%%%%%%%%%%%%%%%%%%%%%%%%%%%%%%%%%%%%%%%%%%%%%%%%%%%%%%%%%%%%%

Die vollständige Matrix $\Sigma = [\sigma_{s,t}]$ integriert alle Segment-Typ-Kombinationen.

\subsection{Die Matrix mit Konfidenzintervallen}

\begin{table}[htbp]
\centering
\caption{Segment-Intervention Effectiveness Matrix $\sigma_{s,t}$ mit 95\% CI}
\label{tab:19-sigma-matrix-full}
\renewcommand{\arraystretch}{1.2}
\footnotesize
\begin{tabular}{@{}lcccccccc@{}}
\toprule
\textbf{Segment} & \textbf{T1} & \textbf{T2} & \textbf{T3} & \textbf{T4} & \textbf{T5} & \textbf{T6} & \textbf{T7} & \textbf{T8} \\
\midrule
Present-Biased & $0.5$ & $0.8$ & \cellcolor{green!25}$1.7$ & $1.4$ & $0.6$ & $0.5$ & $0.8$ & \cellcolor{green!25}$1.9$ \\
Social-Oriented & $0.7$ & $1.5$ & $1.0$ & $0.8$ & \cellcolor{green!25}$1.6$ & \cellcolor{green!25}$1.8$ & $0.6$ & $0.7$ \\
Autonomy-Seeking & $1.3$ & $0.5$ & \cellcolor{green!25}$1.5$ & \cellcolor{red!25}$-0.2$ & $0.9$ & \cellcolor{red!25}$-0.3$ & $0.6$ & $1.4$ \\
Loss-Averse & $0.8$ & $0.9$ & $1.3$ & \cellcolor{green!25}$1.7$ & $0.7$ & $1.1$ & \cellcolor{green!25}$1.6$ & $0.8$ \\
Sophisticates & $0.6$ & $1.0$ & $1.2$ & $1.0$ & $0.8$ & $0.9$ & $0.9$ & \cellcolor{green!25}$1.8$ \\
Naifs & $0.4$ & $0.7$ & \cellcolor{green!25}$1.9$ & $0.6$ & $0.5$ & $0.8$ & $0.7$ & $0.4$ \\
\bottomrule
\end{tabular}
\end{table}

\textbf{Legende:} T1=Information, T2=Feedback, T3=Choice Architecture, T4=Temporal, T5=Identity, T6=Social, T7=Financial, T8=Commitment

\subsection{Nutzung der Matrix für Interventions-Design}

\textbf{Schritt 1:} Segment-Verteilung der Zielgruppe schätzen

\textbf{Schritt 2:} Für jedes Segment den Typ mit maximalem $\sigma$ identifizieren

\textbf{Schritt 3:} Gewichteten erwarteten Effekt berechnen:
\begin{equation}
E_{\text{segment-adjusted}} = \sum_s p_s \cdot E_{\text{base}} \cdot \sigma_{s,t^*(s)}
\label{eq:19-expected-effect}
\end{equation}

\textbf{Schritt 4:} Reaktanz-Risiken für Autonomy-Segment prüfen ($\sigma < 0$?)


%%%%%%%%%%%%%%%%%%%%%%%%%%%%%%%%%%%%%%%%%%%%%%%%%%%%%%%%%%%%%%%%%%%%%%%%%%%%%%%
\section{Targeting-Strategien in der Praxis}
\label{sec:19-targeting}
%%%%%%%%%%%%%%%%%%%%%%%%%%%%%%%%%%%%%%%%%%%%%%%%%%%%%%%%%%%%%%%%%%%%%%%%%%%%%%%

Drei strategische Ansätze für segment-spezifisches Targeting.

\subsection{Strategie 1: Segment-First}

\begin{tcolorbox}[colback=blue!5!white, colframe=blue!75!black,
    title=Strategie: Segment-First Targeting]
\textbf{Ablauf:}
\begin{enumerate}[nosep]
\item Segment-Zugehörigkeit vor Intervention identifizieren
\item Jedes Segment erhält maßgeschneiderte Intervention
\item Maximiert $\sigma$, minimiert Reaktanz-Risiko
\end{enumerate}

\textbf{Pro:} Höchste Effektivität, geringste Reaktanz\\
\textbf{Contra:} Erfordert Segment-Daten vor Intervention
\end{tcolorbox}

\textbf{Segment-Indikatoren für Praxis:}
\begin{itemize}
\item \textbf{Present-Biased:} Historisches Verhalten (Prokrastination), Selbstbericht
\item \textbf{Social-Oriented:} Social Media Nutzung, Gruppen-Mitgliedschaften
\item \textbf{Autonomy-Seeking:} Opt-out-Raten, Beschwerden über ``Bevormundung''
\end{itemize}

\subsection{Strategie 2: Self-Selection Menu}

\begin{tcolorbox}[colback=blue!5!white, colframe=blue!75!black,
    title=Strategie: Self-Selection Menu]
\textbf{Ablauf:}
\begin{enumerate}[nosep]
\item Mehrere Interventions-Optionen anbieten
\item Individuen wählen selbst (nutzt private Information)
\item Respektiert Autonomie (reduziert Reaktanz)
\end{enumerate}

\textbf{Pro:} Nutzt private Information, geringe Reaktanz\\
\textbf{Contra:} Naifs wählen möglicherweise suboptimal
\end{tcolorbox}

\textbf{Menu-Design-Prinzipien:}
\begin{itemize}
\item Commitment-Option für Sophisticates (``Hilf mir, dran zu bleiben'')
\item Social-Option für Social-Oriented (``Gemeinsam mit Kollegen'')
\item Autonome Option für Autonomy-Seeking (``Mein eigener Plan'')
\item Default für Naifs (keine aktive Wahl erforderlich)
\end{itemize}

\subsection{Strategie 3: Adaptive Targeting}

\begin{tcolorbox}[colback=blue!5!white, colframe=blue!75!black,
    title=Strategie: Adaptive Targeting]
\textbf{Ablauf:}
\begin{enumerate}[nosep]
\item Starte mit segmentneutraler Intervention
\item Beobachte Reaktionen, lerne Segment-Zugehörigkeit
\item Personalisiere in Folge-Interventionen
\end{enumerate}

\textbf{Pro:} Kein a-priori Wissen erforderlich\\
\textbf{Contra:} Initiale Effektivitätsverluste, längere Lernphase
\end{tcolorbox}

\textbf{Adaptive Signale:}
\begin{itemize}
\item Opt-out aus Default → Autonomy-Segment
\item Reaktion auf Social Comparison → Social-Segment
\item Nutzung von Commitment-Tools → Sophisticate
\end{itemize}


%%%%%%%%%%%%%%%%%%%%%%%%%%%%%%%%%%%%%%%%%%%%%%%%%%%%%%%%%%%%%%%%%%%%%%%%%%%%%%%
\section{Worked Examples}
\label{sec:19-examples}
%%%%%%%%%%%%%%%%%%%%%%%%%%%%%%%%%%%%%%%%%%%%%%%%%%%%%%%%%%%%%%%%%%%%%%%%%%%%%%%

\subsection{Beispiel 1: Pensionskasse mit 500 Mitarbeitern}

\textbf{Kontext:} HR will Sparquote erhöhen. Budget für eine Kampagne.

\textbf{Segment-Diagnose} (via Survey):
\begin{itemize}[nosep]
\item Present-Biased: 35\% (175 MA)
\item Social-Oriented: 40\% (200 MA)
\item Autonomy-Seeking: 25\% (125 MA)
\end{itemize}

\textbf{Ansatz A: One-Size-Fits-All (Social Norm Message)}

\begin{equation}
E_A = 0.35 \times 0.5 + 0.40 \times 1.8 + 0.25 \times (-0.3) = 0.175 + 0.72 - 0.075 = 0.82
\label{eq:19-example-a}
\end{equation}

\textbf{Problem:} Autonomy-Segment wird geschädigt ($\sigma = -0.3$)!

\textbf{Ansatz B: Segment-Specific Targeting}

\begin{table}[htbp]
\centering
\caption{Segment-spezifisches Targeting: Pensionskasse}
\label{tab:19-example-pension}
\renewcommand{\arraystretch}{1.3}
\small
\begin{tabular}{@{}llccc@{}}
\toprule
\textbf{Segment} & \textbf{Intervention} & \textbf{$\sigma$} & \textbf{Anteil} & \textbf{Beitrag} \\
\midrule
Present-Biased & Auto-Escalation Default & 1.7 & 35\% & 0.595 \\
Social-Oriented & ``85\% sparen bereits...'' & 1.8 & 40\% & 0.720 \\
Autonomy-Seeking & Selbst-Design Sparplan & 1.4 & 25\% & 0.350 \\
\midrule
\textbf{Gesamt} & & & & \textbf{1.665} \\
\bottomrule
\end{tabular}
\end{table}

\textbf{Effizienzgewinn:} $\frac{1.665}{0.82} = 2.03\times$ (+103\% Effektivität)

\subsection{Beispiel 2: Gesundheits-App mit 10,000 Nutzern}

\textbf{Kontext:} App will Bewegung fördern. Kann personalisierte Nachrichten senden.

\textbf{Adaptive Targeting Implementierung:}

\begin{enumerate}
\item \textbf{Woche 1:} Alle erhalten neutrale Information (T1)
\item \textbf{Woche 2:} Segment-Lernen basierend auf Reaktion
   \begin{itemize}[nosep]
   \item Opt-in für Social Features → Social-Oriented
   \item Aktivierung von Reminder-Settings → Present-Biased
   \item Ablehnung von Push-Notifications → Autonomy-Seeking
   \end{itemize}
\item \textbf{Woche 3+:} Personalisierte Interventionen
\end{enumerate}

\textbf{Ergebnis nach 3 Monaten:}
\begin{itemize}
\item Retention Rate: +45\% vs.\ One-Size-Fits-All
\item Aktivitätslevel: +62\% vs.\ Kontrollgruppe
\item Beschwerden: -80\% (weniger wahrgenommene ``Bevormundung'')
\end{itemize}

\subsection{Beispiel 3: Policy-Kampagne mit Reaktanz-Risiko}

\textbf{Kontext:} Regierung will Energiesparen fördern. 30\% der Bevölkerung sind Autonomy-Seeking.

\textbf{Ursprüngliche Kampagne:} ``Jeder Bürger \textit{sollte} seinen Energieverbrauch um 10\% senken.''

\textbf{Problem:} Explizite normative Sprache → $\sigma = -0.3$ für 30\% der Bevölkerung.

\textbf{Modifizierte Kampagne:}
\begin{itemize}
\item Autonomy-Cue hinzufügen: ``Es ist Ihre Entscheidung, wie Sie Energie sparen.''
\item Optionen anbieten: ``Wählen Sie aus 5 Wegen, die zu Ihrem Lebensstil passen.''
\item Default-Architektur: Opt-out statt Opt-in für Energiespar-Programme
\end{itemize}

\textbf{Ergebnis:} $\sigma_{\text{Autonomy}}$ steigt von $-0.3$ auf $+0.8$ durch Reaktanz-Mitigation.


%%%%%%%%%%%%%%%%%%%%%%%%%%%%%%%%%%%%%%%%%%%%%%%%%%%%%%%%%%%%%%%%%%%%%%%%%%%%%%%
\section{Summary}
\label{sec:19-summary}
%%%%%%%%%%%%%%%%%%%%%%%%%%%%%%%%%%%%%%%%%%%%%%%%%%%%%%%%%%%%%%%%%%%%%%%%%%%%%%%

\begin{tcolorbox}[colback=green!10!white, colframe=green!75!black,
    title=Key Points]

\textbf{Core Message:} One-size-fits-all interventions waste resources and can backfire through reactance. Segment-specific targeting doubles effectiveness.

\textbf{Key Takeaways:}
\begin{enumerate}
\item \textbf{Heterogenität ist real:} ATE verbirgt massive Segment-Unterschiede ($\sigma \in [-0.3, 1.9]$)
\item \textbf{Drei Säulen:} Social Preferences ($\alpha$), Time Preferences ($\beta$), Autonomy ($\rho$)
\item \textbf{Present-Biased:} Commitment ($\sigma = 1.9$) und Defaults ($\sigma = 1.7$) optimal
\item \textbf{Social-Oriented:} Social Norms ($\sigma = 1.8$) und Identity ($\sigma = 1.6$) optimal
\item \textbf{Autonomy-Seeking:} Reaktanz-Gefahr! Invisible Defaults + Autonomy-Cues
\item \textbf{Effizienzgewinn:} Segment-spezifisches Targeting verdoppelt Effektivität
\end{enumerate}

\end{tcolorbox}

\begin{tcolorbox}[colback=red!10!white, colframe=red!75!black,
    title=The Golden Rule of Segment Targeting]
\textbf{Niemals} eine explizite normative Intervention bei einem Autonomy-Segment deployen ohne Autonomie-Cue. $\sigma$ kann negativ werden!
\end{tcolorbox}

%------------------------------------------------------------------------------
\paragraph{Formal Details.}
Die vollständige formale Behandlung findet sich in:
\begin{itemize}[nosep]
\item Appendix HHH (METHOD-TOOLKIT) --- F9 (Target Segment), Axiom HHH-5
\item Appendix BA (FORMAL-SEGMENT) --- Segment-Definitionen, Axiome BA-1 bis BA-3
\item Appendix K (LIT-FEHR) --- Fehr's Heterogenitäts-Forschung
\item Appendix BBB (CORE-WHERE) --- Literatur-basierte $\sigma$-Schätzungen
\end{itemize}

%------------------------------------------------------------------------------
\paragraph{What Comes Next.}
\textbf{Kapitel 20} integriert Journey (Kapitel 18) und Segment (Kapitel 19) zu \textbf{kohärenten Interventions-Portfolios}. Die Complementarity Matrix $\gamma_{ij}$ bestimmt, ob Interventionen synergieren oder interferieren --- und wie wir sie optimal kombinieren.


% -----------------------------------------------------------------------------
% READING PATH BOX
% -----------------------------------------------------------------------------

\begin{tcolorbox}[colback=gray!5!white, colframe=gray!75!black,
    title=Empfohlene Lesepfade]

\textbf{Abgeschlossen:} Chapter 19 (Segment-Specific Interventions)

\textbf{Sie verstehen jetzt:}
\begin{itemize}[nosep]
\item Warum ATE für Verhaltensinterventionen irreführend ist
\item Wie der Segment-Multiplier $\sigma_s$ funktioniert
\item Welche Segmente existieren und wie sie reagieren
\item Warum Reaktanz-Management kritisch ist
\end{itemize}

\textbf{Nächste Schritte:}
\begin{itemize}[nosep]
\item \textbf{Kapitel 20}: Intervention Portfolios --- Optimale Kombination von Interventionen
\end{itemize}

\textbf{Für Vertiefung:} Appendix BA (Segment-Axiome), Appendix K (Fehr's Heterogenitäts-Forschung)
\end{tcolorbox}


% =============================================================================
% POLICY IMPLICATIONS
% =============================================================================

\section*{Policy Implications}

\begin{enumerate}
\item \textbf{Für Pensionskassen/HR:} Segment-Diagnose vor Kampagnen-Design. Auto-Enrollment für Present-Biased, Peer-Comparison für Social-Oriented, Self-Design für Autonomy-Seeking. Niemals One-Size-Fits-All.

\item \textbf{Für Gesundheitswesen:} Reaktanz-Risiko bei autonomie-orientierten Patienten ernst nehmen. ``Ärztliche Anordnungen'' können kontraproduktiv sein. Optionen anbieten statt vorschreiben.

\item \textbf{Für Policy-Maker:} Normative Kampagnen (``Bürger sollten...'') riskieren Backlash bei 25--30\% der Bevölkerung. Autonomie-Cues und Optionen reduzieren Reaktanz.

\item \textbf{Für Unternehmen:} Naifs brauchen Defaults, Sophisticates brauchen Commitment-Angebote. Segment-Differenzierung ist kosteneffizienter als universelle Programme.
\end{enumerate}


% =============================================================================
% END OF CHAPTER 19
% =============================================================================
